\documentclass[a4paper,12pt]{article}
\usepackage{graphicx} 
\usepackage{hyperref}
\usepackage{float}
\usepackage{amsmath}
\usepackage{enumitem}
\usepackage{amsfonts}
\usepackage{algorithm}
\usepackage{algorithmic}

% \begin{figure}[H]
%     \centering
%     \includegraphics[width=0.7\textwidth]{filename.png}
%     \caption{Your figure caption here.}
%     \label{fig:yourlabel}
% \end{figure}

\begin{document}

\title{Assignment 5 - Machine Learning \\
Theoretical Questions}
\author{Mohammad Hossein Basouli}
\date{\today}
\maketitle

\section*{Question 1}
\noindent\textbf{Limitations:}
\begin{itemize}
    \item Assumes spherical and convex shapes for each cluster.
    \item Requires specification of number of clusters.
    \item Gives poor results when initialized badly.
\end{itemize}

\noindent\textbf{Initialization Methods:}
\begin{itemize}
    \item \textbf{Random:} Randomly samples \textit{K} distinct data points as centroids.
    \item \textbf{K-Means++:} Selects the first centroid randomly, and then each subsequent centroid is chosen from the remaining points with a probability proportional to its squared distance from the nearest existing centroid.
\end{itemize}

\noindent\textbf{Can K-Means be used with categorical data?}
Not directly, but it's possible if we encode the categorical features using a feature encoding method such as \textit{Label Encoding}.

\section*{Question 2}
It could be due to setting an small $\epsilon$.

\section*{Question 3}
\begin{enumerate}[label=\alph*)]
  \item Because it doesn't assign a given data point to a certain cluster, instead it assign the data point to each of the clusters, with a probability, but hard clustering methods do the opposite.
  \item Yes. Below is the explanation for anomaly detection using each of the methods:
  \begin{itemize}
    \item \textbf{DBSCAN:} At the end of alogithm's execution, the points that have not been assigned to any clusters, are anomalies.
    \item \textbf{GMM:} The data points with a very low likelihood of occurance, after running some iterations of the algorithm, are considered as anomalies.
  \end{itemize}
\end{enumerate}

\section*{Question 4}
\begin{enumerate}[label=\alph*)]
  \item Their difference is written below. Also we can't use \textbf{Silhouette Score} for a dataset which includes only one cluster, since one of parameters in the computation is the average distance from a given point and all other points of the nearest neighbor cluster.  
  \begin{itemize}
    \item \textbf{Silhouette Score:} Computationally inefficient - point-wise calculation of score - higher value is better.
    \item \textbf{Davis-Bouldin Score:} Computationally efficient - cluster-wise calculation of score - lower value is better.
  \end{itemize}

  \item When having well-separated clusters, the silhouette score will be close to 1, and as the clusters tend to overlap with each other, the score decreases to -1. 
\end{enumerate}

\section*{Question 5}
\begin{enumerate}[label=\alph*)]
    \item 
    \begin{itemize}
        \item t-SNE is more computationally expensive than PCA.
        \item t-SNE can capture non-linear patterns but PCA can only capture linear patterns.
        \item t-SNE is primarily used for visualization in 2D or 3D space but PCA could also be used for feature reduction.
    \end{itemize}

    \item Yes that's possible. e.g. we can first apply PCA for feature reduction and then t-SNE for visualization.
    \item \begin{itemize}
        \item Dynamic Time Warping (DTW) is the most appropriate distance metric for clustering time series data because it can compare sequences that vary in speed or alignment.
        \item DTW works by non-linearly warping the time axis of one sequence to optimally align it with another, minimizing the cumulative distance between their matched points.
        \item Unlike Euclidean distance, which assumes fixed one-to-one correspondence between time steps, DTW accounts for time shifts and distortions, making it more robust for real-world temporal data.
        \item As a result, DTW is well-suited for clustering time series where similar patterns may be temporally misaligned but structurally similar.
    \end{itemize} 
\end{enumerate}

\section*{Question 6}
\[
X = 
\begin{bmatrix}
90 & 60 & 90 \\
90 & 90 & 30 \\
60 & 60 & 60 \\
\end{bmatrix}
\]

\subsection*{Step 1: Center the Data}

Compute the mean of each column:

\[
\mu = \frac{1}{3} \sum X = 
\begin{bmatrix}
\frac{90 + 90 + 60}{3} \\
\frac{60 + 90 + 60}{3} \\
\frac{90 + 30 + 60}{3}
\end{bmatrix}
=
\begin{bmatrix}
80 \\
70 \\
60
\end{bmatrix}
\]

Subtract the mean vector from each row:

\[
X_{\text{centered}} = X - \mu^T = 
\begin{bmatrix}
90 - 80 & 60 - 70 & 90 - 60 \\
90 - 80 & 90 - 70 & 30 - 60 \\
60 - 80 & 60 - 70 & 60 - 60 \\
\end{bmatrix}
=
\begin{bmatrix}
10 & -10 & 30 \\
10 & 20 & -30 \\
-20 & -10 & 0 \\
\end{bmatrix}
\]

\subsection*{Step 2: Compute the Covariance Matrix}

\[
\Sigma = \frac{1}{n-1} X_{\text{centered}}^T X_{\text{centered}}
\]

\[
\Sigma = \frac{1}{2}
\begin{bmatrix}
10 & 10 & -20 \\
-10 & 20 & -10 \\
30 & -30 & 0 \\
\end{bmatrix}^T
\begin{bmatrix}
10 & 10 & -20 \\
-10 & 20 & -10 \\
30 & -30 & 0 \\
\end{bmatrix}
= \frac{1}{2}
\begin{bmatrix}
1000 & -500 & 600 \\
-500 & 600 & -900 \\
600 & -900 & 1800 \\
\end{bmatrix}
\]

\[
\Sigma = 
\begin{bmatrix}
500 & -250 & 300 \\
-250 & 300 & -450 \\
300 & -450 & 900 \\
\end{bmatrix}
\]

\subsection*{Step 3: Find Eigenvalues and Eigenvectors}

Solve the characteristic equation:

\[
\det(\Sigma - \lambda I) = 0
\]

This yields eigenvalues:
\[
\lambda_1 = 1205.47,\quad \lambda_2 = 336.20,\quad \lambda_3 = 158.32
\]

Corresponding eigenvectors (normalized):

\[
v_1 =
\begin{bmatrix}
0.35 \\
-0.49 \\
0.80 \\
\end{bmatrix},
\quad
v_2 =
\begin{bmatrix}
0.84 \\
0.54 \\
-0.05 \\
\end{bmatrix}
\]

\subsection*{Step 4: Project the Data onto the First Two Principal Components}

Let \( W = [v_1 \; v_2] \). Then the reduced data:

\[
X_{\text{reduced}} = X_{\text{centered}} \cdot W
\]

\[
X_{\text{reduced}} =
\begin{bmatrix}
10 & -10 & 30 \\
10 & 20 & -30 \\
-20 & -10 & 0 \\
\end{bmatrix}
\cdot
\begin{bmatrix}
0.35 & 0.84 \\
-0.49 & 0.54 \\
0.80 & -0.05 \\
\end{bmatrix}
=
\begin{bmatrix}
30.5 & 0.6 \\
-30.7 & 14.6 \\
0.2 & -15.2 \\
\end{bmatrix}
\]

\subsection*{Conclusion}

The data projected into the 2D space defined by the first two principal components is:

\[
\begin{bmatrix}
30.5 & 0.6 \\
-30.7 & 14.6 \\
0.2 & -15.2 \\
\end{bmatrix}
\]

This concludes the PCA procedure.

\section*{Question 7}
Clustering may benefit \textbf{Linear Regression} in two ways:
\begin{itemize}
    \item Detection of anomalies and removing them before applying \textbf{Linear Regression}.
    \item Applying clustering in order to divide the data into multiple group and applying \textbf{Linear Regression} on each separately.
\end{itemize}

\section*{Question 8}
This shows that each data point is at an equal distance to all of its other neighbors and all are within a single cluster.
\end{document}